% !TEX root = ../main.tex

\begin{figure*}[t]
  \centering
  \begin{subfigure}[t]{0.24\textwidth}
    \centering
    \includegraphics[width=\linewidth]{images/sedd_bars.pdf}\\[-1mm]
    \includegraphics[width=\linewidth]{images/sedd.pdf}
    \caption{SEDD}
    \label{fig: sedd results}
  \end{subfigure}
  \hfill
  \begin{subfigure}[t]{0.24\textwidth}
    \centering
    \includegraphics[width=\linewidth]{images/mdlm_bars.pdf}\\[-1mm]
    \includegraphics[width=\linewidth]{images/mdlm.pdf}
    \caption{MDLM/SDTT}
    \label{fig: mdlm results}
  \end{subfigure}
  \hfill
  \begin{subfigure}[t]{0.24\textwidth}
    \centering
    \includegraphics[width=\linewidth]{images/duo_ancestral_bars.pdf}\\[-1mm]
    \includegraphics[width=\linewidth]{images/duo_ancestral.pdf}
    \caption{Duo Ancestral ($^{\mathrm{a}}$)}
    \label{fig: duo results}
  \end{subfigure}
  \hfill
  \begin{subfigure}[t]{0.24\textwidth}
    \centering
    \includegraphics[width=\linewidth]{images/duo_greedy_bars.pdf}\\[-1mm]
    \includegraphics[width=\linewidth]{images/duo_greedy.pdf}
    \caption{Duo Greedy-Tail ($^{\mathrm{g}}$)}
    \label{fig:duo-greedy-results}
  \end{subfigure}

  \caption{\textbf{Performance comparisons across all benchmarks.} \textit{Top row:} Generation quality/diversity metrics at specific step counts. IDLM-SEDD matches baseline quality with $4\times$ fewer steps ($1024\rightarrow256$), IDLM-MDLM achieves $64\times$ reduction ($1024\rightarrow16$), and IDLM-DCD further reduces to $8$ steps ($128\times$) under Ancestral and $4$ steps ($256\times$) under Greedy-Tail sampling. \textit{Bottom row:} GenPPL efficiency curves. IDLM variants consistently achieve better GenPPL across all step counts while maintaining comparable \textcolor{myorange}{entropy}, with particularly strong improvements in the \underline{low-step} regime.}
  \label{fig:all-results}
  \vspace{-4mm}
\end{figure*}

\vspace{-2mm}
\section{Experiments}
\label{sec:exps}
\vspace{-2mm}
  This section demonstrates the applicability of our IDLM distillation method across various types of DLMs. To this end,
  we conduct experiments using \underline{available checkpoints} of pretrained SEDD, MDLM, and Duo teacher models trained
on OpenWebText (OWT; \citealp{gokaslan2019openwebtext}). For evaluation, we follow the protocol of prior
works~\citep{dieleman2022continuous, lou2024discrete, deschenaux2024beyond, sahoo2024simple, sahoo2025the}
and report GPT-2 Large~\citep{radford2019language} generative perplexity (GenPPL) as well as average sequence entropy to
measure diversity. Comprehensive results are summarized in Table~\ref{tab:all results}. Additional
\underline{technical details} for all evaluated methods are provided in Appendix~\ref{app: experimental details}.
We also provide the \underline{uncurated} text samples in Appendix~\ref{app: qualitative samples}.
All implementations were developed in PyTorch, and the code will be published.


\ecoparagraph{Distillation of SEDD.}
  To the best of our knowledge, no distillation methods have been developed specifically for the SEDD model.
  This section therefore focuses on demonstrating the effectiveness of our distillation approach in reducing the number
  of sampling steps required to achieve comparable generation quality of the teacher model. In the case of SEDD, only
  checkpoints for the absorbing-process variant are available, so we distill this version and refer to it simply as SEDD
  for clarity.
  As shown in Figure~\ref{fig: sedd results}, our method, \texttt{IDLM-SEDD}, consistently outperforms SEDD in terms of
GenPPl across all sampling steps. Notably, IDLM-SEDD maintains high sequence entropy at low sampling steps. IDLM-SEDD
achieves a $4\times$ reduction in sampling steps, accelerating generation from $1024$ to $256$ steps, while preserving
both GenPPL and output diversity.

\ecoparagraph{Distillation of MDLM.}
  This section addresses two primary objectives: (1) demonstrating the feasibility of MDLM model distillation within our
  framework and (2) comparing with the SDTT~\citep{deschenaux2024beyond} a leading method for MDLM distillation.
As shown in Figure~\ref{fig: mdlm results}, our method, \texttt{IDLM-MDLM}, exhibits reduced entropy at large sampling steps, indicating lower diversity in that regime. However, in the low–sampling-step regime, IDLM-MDLM achieves strong performance in terms of both GenPPL and entropy. In particular, IDLM-MDLM attains a $64\times$ acceleration for both MDLM and SDTT, reducing the number of sampling steps from $1024$ to $16$ while preserving comparable generation quality and diversity metrics.

\ecoparagraph{Distillation of Duo.}
  We focus on Duo~\citep{sahoo2025the} and its distillation baseline, Duo-DCD. We distill both the original Duo and
Duo-DCD (treating it as a teacher), denoted \texttt{IDLM-Duo} and \texttt{IDLM-DCD}. Results are in Figure~\ref{fig: duo results}.
With Ancestral sampling, IDLM-Duo matches Duo-DCD with $64\times$ fewer steps ($1024 \to 16$).
IDLM-DCD yields a $128\times$ speed-up ($1024 \to 8$) over the original Duo. With Greedy-Tail sampling,
IDLM-Duo achieves a $64\times$ reduction ($1024 \to 16$), while IDLM-DCD reaches a $256\times$ acceleration ($1024 \to 4$),
all while maintaining comparable metrics.
\vspace{-3mm}
\section{Conclusion}
\vspace{-3mm}
\label{sec:conclusion}
In this work, we introduce \emph{Inverse-distilled Diffusion Language Models} (IDLM), a framework that accelerates
discrete diffusion models by extending Inverse Distillation to the discrete domain.
To ensure valid and stable training, we establish the \emph{theoretical uniqueness} of the objective and propose
\emph{gradient-stable relaxations} using a simplex-parameterized generator.
Empirically, IDLM achieves strong performance across diverse architectures, accelerating \textbf{SEDD by
  $\mathbf{4\times}$},
\textbf{MDLM by $\mathbf{64\times}$}, and \textbf{Duo by up to $\mathbf{256\times}$}, reducing $\mathbf{1024}$ steps to
just $\mathbf{4}$,
while \textit{\textbf{matching}} the teacher's generation quality.
These results position IDLM as a robust solution for practical, low-latency text generation with diffusion models.






% \subsection{Concerns regarding the evaluation protocol used in prior works}
